\section{Background}
\label{sec:background}

\paragraph{Confidential VMs and SEV-SNP.}
A Confidential VM encrypts and integrity-protects an entire guest so that the
host, hypervisor and cloud operator are outside its trust boundary. AMD SEV-SNP
adds strong memory integrity and a hardware-rooted attestation report; Azure
exposes this via Guest Attestation and the Microsoft Azure Attestation (MAA)
service, which returns a signed JWT asserting, among other claims,
\texttt{x-ms-attestation-type=sevsnpvm}~\cite{misono2024cvms,paradzik2024sevsnp}.
Remote attestation lets a relying party verify such evidence; the IETF RATS
architecture standardises the Attester/Verifier/Relying-Party roles and insists
that appraisal be independent of evidence
production~\cite{birkholz2023rats,chen2019opera}.

\paragraph{Kubernetes scheduling and admission.}
The default scheduler binds pods to nodes using the scheduling framework
(PreFilter/Filter/Score/Reserve/Permit/Bind); a cluster may also run additional
schedulers for pods that request them by name. Placement inputs --- labels,
\texttt{nodeSelector}, node affinity, \texttt{RuntimeClass} --- are ordinary,
attacker-writable cluster objects with no cryptographic backing. Admission
webhooks and \texttt{schedulingGates} can gate a pod until an external controller
approves it, but the gate check and the eventual bind are separate events, so a
check-then-ungate design has an inherent TOCTOU gap. \texttt{RuntimeClass}
selects a runtime handler (e.g.\ a confidential-container runtime) but carries no
freshness, revocation or node-identity semantics of its own.

\paragraph{Confidential containers vs.\ confidential nodes.}
Confidential Containers (CoCo) with Trustee run each pod inside its own TEE and
gate secret release on pod-level attestation; recent work demonstrates pod-level
remote attestation on Kubernetes~\cite{yang2026dstack,coco2024,k8s2023confidential}.
On confidential-VM node pools, by contrast, the TEE is the \emph{node}: evidence
attests the confidential worker VM, not each pod. Our work targets the latter and
is explicit about it; it is complementary to pod-level systems.
